\documentclass[12pt,a4paper,oneside]{report}

\usepackage[margin=1in]{geometry}
\usepackage[T1]{fontenc}
\usepackage[utf8]{inputenc}
\usepackage{lmodern}
\usepackage{microtype}
\usepackage{booktabs}
\usepackage{array}
\usepackage{longtable}
\usepackage{enumitem}
\usepackage{xcolor}
\usepackage{hyperref}
\usepackage{listings}
\usepackage{amsmath}
\usepackage{amssymb}
\usepackage{graphicx}
\usepackage{float}

\hypersetup{
  colorlinks=true,
  linkcolor=blue!45!black,
  citecolor=blue!45!black,
  urlcolor=blue!45!black,
  pdftitle={Accelera AutoML: Design and Implementation},
  pdfauthor={Individual Work Draft}
}

\lstdefinestyle{pythonstyle}{
  language=Python,
  basicstyle=\ttfamily\small,
  keywordstyle=\color{blue!60!black},
  commentstyle=\color{green!40!black},
  stringstyle=\color{red!55!black},
  frame=single,
  breaklines=true,
  columns=fullflexible,
  keepspaces=true,
  showstringspaces=false
}

\newcommand{\systemname}{Accelera AutoML}
\newcommand{\modulepath}[1]{\texttt{#1}}
\newcommand{\todoitem}[1]{\textcolor{red!65!black}{[TODO: #1]}}

\title{
  \vspace{-1.5cm}
  \textbf{Accelera AutoML}\[0.4em]
  \Large Automated Model Selection, Resource-Aware Search,\\
  Meta-Learning, and Ensemble Construction
}
\author{Individual Work Draft\\[0.5em]\todoitem{Add student name and identifier}}
\date{June 2026}

\begin{document}

\maketitle

\begin{abstract}
This work addresses automated model selection for supervised tabular data in
the Accelera framework. The implemented module exposes classifier and regressor
interfaces and searches conditional model spaces under explicit limits on the
number and duration of trials. Candidate configurations are evaluated by
cross-validation at three fidelity levels. After an initial phase, a
random-forest surrogate ranks sampled configurations by expected improvement.
The search can also start from configurations retrieved using dataset
metafeatures. At the end of a run, the best individual configuration is
compared with an optional voting or stacked ensemble.

The contribution of this draft is the design and implementation of the module,
not a claim that it outperforms existing AutoML packages. Such a claim requires
the controlled experiments that will be added in the later evaluation chapters.
The present chapters therefore describe the problem, relevant literature,
requirements, method, source-code organization, and known implementation
constraints. Each implementation claim is tied to a class, parameter, or data
flow that exists in the repository.
\end{abstract}

\tableofcontents
\listoftables
\clearpage

% -----------------------------------------------------------------------------
\chapter{Introduction}
\label{chap:introduction}

\section{Background}

Machine-learning development is rarely limited to fitting one estimator. A
practical workflow requires the practitioner to select a model family, define
its hyperparameters, decide whether features require scaling, choose an
evaluation metric, allocate computational resources, and compare alternatives
without leaking validation information. These decisions interact. For
example, a support-vector machine may benefit from feature scaling, while a
tree-based model usually does not require it. A configuration that performs
well on a small sample may fail when trained at full fidelity, and a highly
accurate model may be unsuitable if its search cost exceeds the available
budget.

Automated machine learning (AutoML) treats these choices as a structured
optimization problem. Instead of manually testing a short list of models, an
AutoML system defines a search space and evaluates configurations under a
consistent protocol. The system must balance exploration of new regions with
exploitation of observations collected during the current search. It must also
handle practical failures such as invalid parameter combinations, estimator
exceptions, and trials that exceed their time limit.

\section{Problem Statement}

The central problem addressed by this work is:

\begin{quote}
How can Accelera search classification and regression configurations under
explicit trial and time constraints, avoid spending full resources on every
candidate, and return the best model observed under a consistent evaluation
protocol?
\end{quote}

This problem contains several subproblems:

\begin{enumerate}[leftmargin=*]
  \item representing heterogeneous and conditional hyperparameter spaces;
  \item comparing configurations using a fair task-specific evaluation
        protocol;
  \item allocating low and high fidelities to candidates efficiently;
  \item transferring useful configurations from related historical datasets;
  \item selecting a final individual model or ensemble; and
  \item containing failures so that one estimator does not terminate the
        complete search.
\end{enumerate}

\section{Motivation}

Manual model selection is difficult to reproduce because it often depends on
informal trial-and-error decisions. It is also inefficient when many model
families have large hyperparameter spaces. A reusable AutoML component provides
three practical benefits. First, all candidates are evaluated by the same
metric and validation policy. Second, resource limits become explicit through
parameters such as \texttt{time\_budget}, \texttt{n\_trials}, and
\texttt{per\_run\_time\_limit}. Third, the complete run history can be exposed
as a leaderboard for inspection rather than returning an unexplained model.

The module is implemented inside Accelera rather than as an isolated benchmark
script. This decision makes installation, generated public API files, packaged
meta-learning data, and future integration with the preprocessing pipeline part
of the same project. The public classes follow scikit-learn's estimator pattern:
configuration is supplied to the constructor, training occurs in
\texttt{fit}, and the fitted object provides \texttt{predict}, \texttt{score},
and access to its leaderboard and run history.

\section{Development Context}

The implementation was developed as a separate module under
\modulepath{accelera/src/accelera\_automl}. This distinction matters because
Accelera already contained an automatic-preprocessing module. The new module is
responsible for choosing and fitting estimators; it does not replace the
existing preprocessing reports and dataset preparation workflow.

Several design choices came directly from problems observed while integrating
the module. Trials need process isolation because an estimator that exceeds its
budget cannot always be stopped safely from a worker thread. Search parallelism
also has to be separated from estimator parallelism to avoid creating more
workers than the machine can support. Meta-learning JSON files must be resolved
relative to the installed package instead of a developer's local path. Finally,
the public exports have to be generated from the source tree so that an
installed wheel and a source checkout expose the same import path. These issues
are implementation concerns, but they influence the proposed method because
they determine which resource controls and failure states are visible to the
optimizer.

\section{Aim and Objectives}

The aim of this work is to design and implement a resource-aware AutoML engine
for tabular classification and regression. Its objectives are:

\begin{itemize}[leftmargin=*]
  \item provide \texttt{AutoMLClassifier} and \texttt{AutoMLRegressor};
  \item search multiple model families and their conditional parameters;
  \item support user-defined scoring, cross-validation, and reproducible seeds;
  \item evaluate candidates at progressively higher fidelity;
  \item use observed trials to guide later sampling;
  \item warm-start search using dataset metafeatures and stored historical runs;
  \item construct voting or stacked ensembles from competitive candidates;
  \item support parallel trials and per-trial timeout isolation; and
  \item return the final model, best score, run history, and leaderboard.
\end{itemize}

\section{Contributions}

The current implementation makes the following engineering and methodological
contributions:

\begin{enumerate}[leftmargin=*]
  \item A shared engine for classification and regression with task-specific
        evaluators and component registries.
  \item A three-stage fidelity schedule that varies sample fraction,
        cross-validation folds, and estimator budget.
  \item A sequential model-based optimizer using a random-forest surrogate and
        expected improvement over a sampled candidate pool.
  \item A promotion mechanism that advances successful low-fidelity candidates
        according to their relative cost.
  \item Classification and regression metafeatures for retrieving warm-start
        configurations from packaged metadata.
  \item Voting and stacked ensemble construction, including forward selection
        of diverse base models.
  \item A stable public namespace:
        \modulepath{accelera.accelera\_automl}.
\end{enumerate}

\section{Scope}

This thesis focuses on supervised tabular classification and regression. The
implemented system assumes data that can be validated by scikit-learn's array
validation utilities. The work does not currently claim support for automated
deep neural architecture search, unsupervised learning, time-series-specific
validation, streaming updates, or distributed cluster execution. These are
possible extensions but are outside the present implementation scope.

\section{Organization of This Draft}

Chapter~\ref{chap:background} introduces the technical foundations and related
AutoML systems. Chapter~\ref{chap:requirements} defines system requirements and
design constraints. Chapter~\ref{chap:method} presents the proposed search,
meta-learning, and ensemble methods. Chapter~\ref{chap:implementation} maps
those methods to the source-code architecture. Chapters concerning experiments,
results, discussion, and conclusion are reserved for later work and are not
included in this file.

% -----------------------------------------------------------------------------
\chapter{Background and Related Work}
\label{chap:background}

\section{Combined Algorithm Selection and Hyperparameter Optimization}

Given a set of algorithms $\mathcal{A}$ and a hyperparameter domain
$\Lambda_a$ for every algorithm $a$, combined algorithm selection and
hyperparameter optimization seeks

\begin{equation}
  (a^*, \lambda^*) =
  \arg\min_{a \in \mathcal{A},\,\lambda \in \Lambda_a}
  \widehat{L}(a, \lambda; D),
\end{equation}

where $\widehat{L}$ is an estimated validation loss on dataset $D$. Auto-WEKA
formalized this combined view by searching over both learning algorithms and
their hyperparameters \cite{autoweka}. The domain is conditional: parameters
belonging to a random forest are inactive when the selected model is logistic
regression. This conditional structure motivates the use of the ConfigSpace
representation in \systemname{}.

\section{Search Strategies}

Grid search evaluates a fixed Cartesian product and scales poorly as the
number of dimensions increases. Random search allocates trials more effectively
when only a subset of hyperparameters strongly affects performance
\cite{randomsearch}. Sequential model-based optimization improves on uninformed
sampling by fitting a surrogate to observed configuration costs. SMAC uses a
random-forest model for algorithm configuration problems and selects candidates
through an acquisition function \cite{smac}.

\systemname{} follows this general model-based pattern. It first evaluates
initial and warm-start configurations, represents each configuration as a
numeric vector, augments that vector with fidelity variables, and fits a
random-forest regressor. Predictions from individual trees provide an empirical
mean and standard deviation. Expected improvement then ranks a randomly sampled
candidate pool. The implementation is inspired by SMAC principles but is a
project-specific optimizer rather than a dependency on the external SMAC
package.

\section{Multi-Fidelity Evaluation}

Full cross-validation on an entire dataset is expensive, particularly for
large ensembles or boosting models. Multi-fidelity optimization first evaluates
candidates using cheaper approximations. Fidelity may be varied by the number
of training instances, the number of validation folds, training iterations, or
model capacity. Poor candidates can be discarded early, while promising ones
receive more resources.

In this work, fidelity is represented by a tuple

\begin{equation}
  f = (s, q, k, b),
\end{equation}

where $s$ is the stage index, $q$ is the sample fraction, $k$ is the number of
cross-validation folds, and $b$ is a multiplicative model budget. The fidelity
vector is included in the surrogate input so that observations from different
stages remain distinguishable.

\section{Meta-Learning for Warm Starts}

Meta-learning attempts to transfer knowledge across machine-learning tasks.
An AutoML system can characterize a new dataset using metafeatures, compare it
with previously evaluated datasets, and schedule successful historical
configurations early. Auto-sklearn combines Bayesian optimization,
meta-learning, and ensemble construction \cite{autosklearn}. The intended
benefit is not to guarantee that a historical configuration remains optimal,
but to reduce the cold-start period before useful observations are available.

\systemname{} computes task-dependent statistical descriptors. Shared
descriptors include the number of instances and features, numeric/categorical
ratios, missing-value ratios, skewness, kurtosis, and categorical symbol
statistics. Classification additionally records class entropy, class-count
statistics, and class probabilities. Regression records target mean, standard
deviation, skewness, and kurtosis. Stored datasets are ranked by distance in
this metafeature space, and candidate configurations from the closest datasets
are converted back into active ConfigSpace objects.

\section{Ensemble Learning}

Ensembles combine several fitted models to reduce the risk of selecting one
unstable estimator. Voting uses the predictions or probabilities of multiple
base estimators. Stacked generalization instead trains a meta-model on
out-of-fold predictions from base models \cite{stacking}. Out-of-fold
generation is essential: fitting the meta-model on predictions made on each
base model's own training observations would produce leakage and an overly
optimistic result.

The Accelera implementation provides both strategies. Stacking uses bagged
base estimators, creates out-of-fold prediction blocks, performs forward
selection, and fits a logistic-regression meta-model for classification or a
ridge-regression meta-model for regression. The original feature matrix may
optionally be appended to the prediction blocks.

\section{Related AutoML Systems}

The systems in Table~\ref{tab:related-systems} were not selected as a general
list of popular AutoML libraries. Each one provides a comparison for a specific
part of the Accelera design. Auto-WEKA establishes the combined algorithm and
hyperparameter-selection problem. Auto-sklearn is the closest conceptual
comparison because it combines configuration optimization, metafeatures, and
ensembles. FLAML is relevant to the use of explicit resource budgets.
AutoGluon Tabular is included for its emphasis on stacking, while TPOT is a
contrasting case because it searches pipeline structures using evolutionary
operators rather than the surrogate-guided method used here. These packages are
related work and possible experimental baselines; their implementations are not
called by \systemname{}.

\begin{table}[H]
\centering
\caption{Selected AutoML systems and their primary design emphasis}
\label{tab:related-systems}
\begin{tabular}{p{0.18\textwidth}p{0.72\textwidth}}
\toprule
System & Design emphasis \\
\midrule
Auto-WEKA & Combined algorithm selection and hyperparameter optimization over
             WEKA components \cite{autoweka}. \\
auto-sklearn & Bayesian optimization, dataset metafeatures, and post-hoc
                ensembles for scikit-learn workflows \cite{autosklearn}. \\
TPOT & Evolutionary search over tree-structured machine-learning pipelines
       \cite{tpot}. \\
FLAML & Cost-aware, lightweight AutoML intended to find strong configurations
        under constrained resources \cite{flaml}. \\
AutoGluon Tabular & Multi-layer stacking and ensembling with a high-level
                     tabular prediction interface \cite{autogluon}. \\
\systemname{} & Conditional model spaces, staged fidelity, surrogate-guided
                 sampling, packaged warm starts, and selectable voting or
                 stacking inside Accelera. \\
\bottomrule
\end{tabular}
\end{table}

The comparison should not be read as an equivalence claim. For example, the
current Accelera optimizer is implemented locally and is smaller in scope than
the mature optimization infrastructure used by established frameworks. The
purpose of the comparison is to identify which ideas are shared, which search
mechanisms differ, and which systems should be used as baselines in the later
experimental chapter.

% -----------------------------------------------------------------------------
\chapter{Requirements and System Analysis}
\label{chap:requirements}

\section{Stakeholders and Use Cases}

The primary user is a data scientist or developer who has a supervised dataset
and wants a competitive baseline without manually enumerating estimators. A
secondary user is a researcher who needs access to the run history,
leaderboard, fidelity stages, and ensemble decisions. A maintainer must be able
to add a model family without rewriting the optimization engine.

The main use cases are:

\begin{itemize}[leftmargin=*]
  \item fit a classifier and obtain labels or probabilities;
  \item fit a regressor and obtain numeric predictions;
  \item restrict search to an allowed subset of model families;
  \item select a metric and cross-validation fold count;
  \item limit the complete search or each individual trial;
  \item enable or disable meta-learning and ensembles;
  \item inspect the best score, leaderboard, and complete run history; and
  \item reproduce a run using a fixed random state.
\end{itemize}

\section{Functional Requirements}

\begin{longtable}{p{0.12\textwidth}p{0.78\textwidth}}
\caption{Functional requirements}\label{tab:functional-requirements}\\
\toprule
ID & Requirement \\
\midrule
\endfirsthead
\toprule
ID & Requirement \\
\midrule
\endhead
FR-01 & The system shall support supervised classification and regression. \\
FR-02 & The public estimators shall provide \texttt{fit}, \texttt{predict}, and
        \texttt{score}; classification shall also provide
        \texttt{predict\_proba} when supported by the selected final model. \\
FR-03 & The system shall validate input arrays before fitting or prediction. \\
FR-04 & The engine shall construct a task-specific conditional configuration
        space from registered components. \\
FR-05 & The evaluator shall compare configurations using cross-validation and
        the selected scoring function. \\
FR-06 & The optimizer shall stop when the trial limit or overall time budget is
        reached. \\
FR-07 & Trials shall support staged fidelity and promotion to higher stages. \\
FR-08 & The search shall optionally begin with meta-learning warm starts. \\
FR-09 & The engine shall optionally construct voting or stacked ensembles. \\
FR-10 & The final result shall include a fitted final model, score,
        leaderboard, ensemble reference, and run history. \\
FR-11 & Trial exceptions and timeouts shall be recorded instead of terminating
        the entire optimization run. \\
FR-12 & Independent trials shall support process-level parallel execution. \\
\bottomrule
\end{longtable}

\section{Non-Functional Requirements}

\begin{itemize}[leftmargin=*]
  \item \textbf{Usability:} API behavior should resemble scikit-learn
        estimators and use explicit keyword parameters.
  \item \textbf{Reproducibility:} stochastic estimators, data splitting, and
        search sampling should receive the configured random state where
        supported.
  \item \textbf{Extensibility:} component definitions, evaluator behavior,
        optimization, meta-learning, and ensembles should remain separate.
  \item \textbf{Fault isolation:} long-running trials should be terminable
        without corrupting the parent search process.
  \item \textbf{Resource control:} nested parallelism should be configurable
        through search, stack, and inner job counts.
  \item \textbf{Inspectability:} failed and successful trials should remain in
        the run history with status, duration, score, cost, and fidelity data.
  \item \textbf{Packaging:} runtime metadata required by meta-learning should
        be distributed with the Python package.
\end{itemize}

\section{Constraints and Assumptions}

The current implementation is constrained by the interfaces of the underlying
estimators. Some model parameters represent iterations, estimators, or depth,
so the meaning of a model-budget multiplier is component-specific. Process
timeouts depend on multiprocessing behavior and may have platform-specific
overhead. In addition, model libraries such as CatBoost, LightGBM, and XGBoost
must be installed for their registered components to be imported.

The evaluator assumes independent and identically distributed tabular samples.
Classification uses stratified splitting where possible, while regression uses
ordinary K-fold or shuffle splitting. Specialized grouped, temporal, or spatial
validation is outside the current design.

\section{Risk Analysis}

\begin{table}[H]
\centering
\caption{Principal technical risks and current mitigations}
\begin{tabular}{p{0.25\textwidth}p{0.30\textwidth}p{0.34\textwidth}}
\toprule
Risk & Effect & Mitigation \\
\midrule
Invalid configuration & Trial exception or meaningless score & Conditional
space definitions and failure results. \\
Long-running estimator & Search exceeds its budget & Per-run process timeout
and termination. \\
Nested parallelism & CPU oversubscription & Separate search, stack, and inner
job controls. \\
Class imbalance & Misleading accuracy or failed folds & Stratification,
class-weight support, and fold-count adjustment. \\
Small datasets & Unstable high-fold validation & Effective fold count bounded
by available class samples. \\
Meta-data mismatch & Invalid historical configuration & Filtering by allowed
models and guarded conversion to ConfigSpace. \\
Overfitted ensemble & Inflated training performance & Out-of-fold predictions
for meta-model fitting. \\
\bottomrule
\end{tabular}
\end{table}

% -----------------------------------------------------------------------------
\chapter{Proposed AutoML Method}
\label{chap:method}

\section{Method Overview}

The complete method is organized as the following sequence:

\begin{enumerate}[leftmargin=*]
  \item validate the feature matrix and target;
  \item resolve compatible or user-allowed model families;
  \item construct the conditional configuration space;
  \item compute dataset metafeatures and retrieve warm-start configurations;
  \item evaluate initial configurations at low fidelity;
  \item fit a surrogate after sufficient finite observations exist;
  \item select new candidates by expected improvement;
  \item promote strong candidates through higher fidelity stages;
  \item fit the best full-fidelity individual configuration;
  \item optionally construct and evaluate an ensemble; and
  \item return whichever fitted candidate has the stronger validation score.
\end{enumerate}

\section{Configuration-Space Construction}

Classification and regression use independent component registries. Each
component specifies a model name, hyperparameter factory, conditional rules,
and estimator builder. The global configuration contains a categorical
\texttt{model\_name}. Model-specific hyperparameters are activated only when
their model is selected. This prevents the optimizer from treating irrelevant
parameters as active decisions.

Before construction, the engine resolves an allowed model set. The user may
provide this set explicitly. The engine can also remove incompatible models
based on dataset properties, such as algorithms whose assumptions conflict
with negative feature values. Restricting the search space reduces invalid
trials and allows domain knowledge to be expressed without modifying the
optimizer.

\section{Evaluation and Cost Conversion}

For configuration $\lambda$, the evaluator builds a model and a list of
candidate preprocessing strategies. Each model--preprocessor pair is scored by
cross-validation, and the best successful preprocessing result is retained.
Let $m(\lambda)$ be the mean validation metric. The optimizer minimizes a cost
$c(\lambda)$ obtained from the metric. For maximization metrics, the evaluator
uses an inverse or negated form so that smaller cost remains better. Failed
evaluations receive a defined failure cost and an error description.

Classification uses stratified folds and stratified subsampling. Regression
uses K-fold validation and ordinary random subsampling. The classifier adjusts
the requested fold count when the smallest class cannot support that number of
folds.

\section{Three-Stage Fidelity Schedule}

The default fidelity schedule is shown in
Table~\ref{tab:fidelity-schedule}.

\begin{table}[H]
\centering
\caption{Default evaluation fidelity stages}
\label{tab:fidelity-schedule}
\begin{tabular}{ccccc}
\toprule
Stage & Sample fraction & CV folds & Model budget & Purpose \\
\midrule
0 & 0.20 & 2 & 0.25 & inexpensive screening \\
1 & 0.50 & 3 & 0.60 & intermediate confirmation \\
2 & 1.00 & configured value & 1.00 & full-fidelity ranking \\
\bottomrule
\end{tabular}
\end{table}

The evaluator scales supported capacity parameters using the model budget. A
configuration is initially scheduled at stage zero. After a successful result,
the optimizer records the cost in a per-configuration state. Once enough
observations exist, only candidates inside the configured top promotion
quantile advance. This design spends the most expensive evaluations on
candidates that were competitive at a cheaper stage.

\section{Surrogate-Guided Candidate Selection}

Every ConfigSpace configuration is converted into a numerical vector. Missing
or inactive dimensions are replaced by finite sentinel values. The vector is
then concatenated with the fidelity representation

\begin{equation}
  x' = [x_{\lambda}, s, q, k, b].
\end{equation}

After the initial exploration phase, a random-forest regressor is fitted to
finite observed costs. If tree $t$ predicts $\hat{c}_t(x)$, the surrogate mean
and standard deviation are estimated as

\begin{align}
  \mu(x) &= \frac{1}{T}\sum_{t=1}^{T}\hat{c}_t(x), \\
  \sigma(x) &= \sqrt{\frac{1}{T}\sum_{t=1}^{T}
  \left(\hat{c}_t(x)-\mu(x)\right)^2}.
\end{align}

For minimization, expected improvement over the best observed cost $c^+$ is

\begin{align}
  I(x) &= c^+ - \mu(x) - \xi, \\
  z(x) &= \frac{I(x)}{\sigma(x)}, \\
  \operatorname{EI}(x) &= I(x)\Phi(z) + \sigma(x)\phi(z),
\end{align}

where $\Phi$ and $\phi$ are the standard normal cumulative distribution and
density functions. The optimizer samples a candidate pool, removes duplicate
vectors, computes EI, and schedules the highest-ranked unseen configuration.
This procedure avoids exhaustive enumeration of a mixed conditional space.

\section{Warm-Start Selection}

For a new dataset, the engine computes the task-specific metafeature vector
$m_D$. Historical datasets are represented by stored vectors $m_i$. A
normalized distance ranks historical datasets by similarity. Configurations
are gathered from the nearest datasets, filtered to the active model set,
deduplicated by model and parameter signature, and converted to ConfigSpace
configurations. The number of datasets, configurations per dataset, and total
warm starts are bounded by user parameters.

Warm starts enter the initial configuration queue before uninformed samples.
They do not bypass evaluation and therefore cannot directly force the final
choice. If a historical configuration is invalid under the current search
space, conversion is skipped safely.

\section{Parallel Execution and Timeout Control}

The optimizer can suggest a batch containing promotions, warm starts, and new
surrogate-selected or random configurations. For parallel batches, each trial
runs in its own process and communicates an \texttt{EvaluationResult} through a
queue. The parent process periodically checks queues and elapsed time. A process
that exceeds \texttt{per\_run\_time\_limit} is terminated and recorded with
timeout status.

Process-level execution is used because a timed-out estimator cannot always be
interrupted reliably from a thread. The design also separates
\texttt{search\_n\_parallel} from internal estimator job counts, allowing the
caller to limit nested parallelism.

\section{Final Model and Ensemble Selection}

The strongest full-fidelity configuration is rebuilt and fitted on all
available training data. If ensembles are disabled, this model is returned
directly. Otherwise, successful run-history candidates are considered for
voting or stacking.

For voting, the engine selects a bounded set of competitive estimators and
constructs a task-specific voting estimator. For stacking, candidate diversity
is considered before selecting base estimators. Every base estimator is wrapped
in bagging and produces out-of-fold predictions. Forward selection begins with
the strongest single base model and adds candidates while they improve the
meta-model score or until the minimum base-model requirement is satisfied.

For classification, probability blocks are aligned across class labels before
being passed to logistic regression. For regression, each base model contributes
one numeric prediction column and ridge regression acts as the meta-model. The
ensemble is evaluated using the same task metric. It replaces the individual
model only when its validation score is at least as strong.

% -----------------------------------------------------------------------------
\chapter{System Design and Implementation}
\label{chap:implementation}

\section{Package Architecture}

The implementation is stored under
\modulepath{accelera/src/accelera\_automl/}. Public exports are generated under
\modulepath{accelera/api/accelera\_automl/}, allowing users to import
\texttt{AutoMLClassifier} and \texttt{AutoMLRegressor} from
\modulepath{accelera.accelera\_automl}.

\begin{table}[H]
\centering
\caption{Main implementation units}
\begin{tabular}{p{0.34\textwidth}p{0.56\textwidth}}
\toprule
File or directory & Responsibility \\
\midrule
\modulepath{base.py} & shared estimator lifecycle, validation, fitted state,
leaderboard, and run-history access \\
\modulepath{classification.py} & classification defaults and engine
construction \\
\modulepath{regression.py} & regression defaults and engine construction \\
\modulepath{core/automl.py} & orchestration, final model fitting, compatibility
checks, and ensemble selection \\
\modulepath{components/} & model registries, hyperparameters, conditions, and
estimator factories \\
\modulepath{configspace\_search\_space.py} & task-specific ConfigSpace
assembly and configuration conversion \\
\modulepath{evaluation.py} & fidelity-aware model construction,
preprocessing comparison, cross-validation, and cost conversion \\
\modulepath{optimization/smac\_optimizer.py} & trial scheduling, surrogate,
expected improvement, promotions, parallel workers, and timeouts \\
\modulepath{meta\_learning/} & metafeature extraction, similarity ranking, and
warm-start retrieval \\
\modulepath{stacked\_ensemble.py} & classification adapters, bagging, forward
selection, and stacked classification \\
\modulepath{stacked\_ensemble\_regression.py} & regression stacking and
forward selection \\
\modulepath{meta\_learning\_data/} & packaged historical classification and
regression metadata \\
\bottomrule
\end{tabular}
\end{table}

\section{Traceability from Method to Source Code}

Table~\ref{tab:traceability} connects the principal claims in
Chapter~\ref{chap:method} to concrete implementation points. This table is
included to distinguish implemented behavior from work that is only proposed
for the future.

\begin{table}[H]
\centering
\caption{Method-to-implementation traceability}
\label{tab:traceability}
\begin{tabular}{p{0.27\textwidth}p{0.31\textwidth}p{0.31\textwidth}}
\toprule
Method element & Implementation point & Observable state or output \\
\midrule
Three fidelity stages &
\texttt{AutoMLEngine.DEFAULT\_FIDELITY\_STAGES} &
stage, sample fraction, folds, and model budget in each trial result \\
Expected improvement &
\texttt{Optimizer.expected\_improvement} &
ranking of the sampled candidate pool \\
Candidate promotion &
\texttt{Optimizer.record\_observation} and
\texttt{should\_promote} &
per-configuration successful stages and promotion queue \\
Meta-learning warm start &
\texttt{get\_meta\_learning\_warmstarts} &
initial ConfigSpace configurations \\
Per-trial timeout &
\texttt{evaluate\_batch\_with\_processes} &
timeout status, duration, and error text \\
Dataset-aware filtering &
\texttt{AutoMLEngine.resolve\_allowed\_models} &
model names removed before space construction \\
Stacked forward selection &
\texttt{forward\_select\_base\_models} &
selected base names and forward-selection score \\
Final comparison &
\texttt{AutoMLEngine.search} &
individual model, optional ensemble, final model, and leaderboard \\
\bottomrule
\end{tabular}
\end{table}

\section{Public Estimator Lifecycle}

\texttt{BaseAutoML} extends scikit-learn's \texttt{BaseEstimator}. Its
\texttt{fit} method validates $X$ and $y$, resets previously fitted state,
builds a task-specific engine, runs the search, and copies the result into
public fitted attributes. Prediction validates $X$ and delegates to the final
model. Calling prediction before fitting raises an explicit runtime error.

The result state includes:

\begin{itemize}[leftmargin=*]
  \item \texttt{best\_model}: strongest fitted individual candidate;
  \item \texttt{ensemble}: fitted ensemble when constructed;
  \item \texttt{final\_model}: model selected for prediction;
  \item \texttt{best\_score}: validation score of the final choice;
  \item \texttt{leaderboard}: ranked model and ensemble summary; and
  \item \texttt{runhistory}: detailed trial records.
\end{itemize}

\section{Engine Collaboration}

The engine uses composition rather than embedding all behavior in one class.
It constructs a configuration space, evaluator, optimizer, and optional
ensemble. This separation allows the evaluator to remain task-aware while the
optimizer operates mainly on configurations, numeric vectors, and scalar costs.

The runtime interaction is:

\begin{verbatim}
User estimator
  -> BaseAutoML.fit
  -> AutoMLEngine.search
     -> ConfigSpace construction
     -> metafeature extraction and warm starts
     -> task evaluator
     -> Optimizer.optimize
        -> staged trials / process workers
        -> surrogate and expected improvement
        -> promotion queue
     -> final individual model
     -> optional voting or stacked ensemble
  -> AutoMLResult
  -> fitted public estimator
\end{verbatim}

\section{Model Component Abstraction}

Classification and regression components encapsulate four related concerns:
their public name, hyperparameter definitions, conditional activation rules,
and estimator construction. This abstraction prevents the optimizer from
containing model-specific \texttt{if} statements. Adding an estimator requires
registering a component and ensuring its parameters can be scaled at the
defined fidelity levels when applicable.

The current registries include estimators from scikit-learn and optional
gradient-boosting libraries. Because all models do not expose identical
interfaces, the component builders normalize values such as textual
representations of \texttt{None}, class weights, and task-specific objectives.

\section{Evaluation Result Model}

Every evaluation returns a structured record rather than only a score. The
record contains the model name, resolved parameters, preprocessing choice,
metric score, optimizer cost, duration, status, error text, stage, sample
fraction, fold count, and model budget. This representation serves three
purposes: it supports the optimizer, preserves failure diagnostics, and enables
leaderboard or experimental analysis without repeating the run.

\section{Meta-Learning Data Packaging}

Classification and regression metadata are stored in separate JSON
directories. File selection depends on task, metric, and classification target
shape. Paths are resolved relative to the installed module rather than a
developer-specific absolute directory. The package configuration explicitly
includes these JSON resources in built wheels so warm starts remain available
after installation.

\section{Ensemble Implementation Details}

The classification stacker includes an adapter that guarantees a probability
interface. If a base model lacks \texttt{predict\_proba}, decision scores are
converted to probabilities where possible; otherwise predicted classes are
encoded into probability-like output. Class columns are aligned so each base
block has consistent semantics.

Both stackers generate out-of-fold predictions using cross-validation and fit
fresh bagged models on all training data for later inference. Forward-selection
state is stored as a fitted attribute together with its score. This explicit
state is also used by verbose structural logging and avoids recomputing the
selection result.

Two implementation rules limit ensemble cost. First, the number of candidate
base models is bounded by \texttt{stacked\_base\_size} or
\texttt{ensemble\_voting\_size}. Second, a requested stacked ensemble is
changed to voting when the dataset contains more than 50,000 rows. This
threshold is a project heuristic rather than a theoretically optimal boundary;
its effect must be evaluated in the later experiments.

\section{Implementation Decisions and Trade-offs}

The optimizer uses a sampled candidate pool instead of optimizing expected
improvement continuously. This is simpler for a conditional space containing
categorical and inactive parameters, but its quality depends on
\texttt{candidate\_pool\_size}. A small pool is cheaper and may miss a useful
region; a large pool increases acquisition overhead.

The default per-run timeout is 60 seconds when no overall budget is supplied.
When \texttt{time\_budget} is available, the default is one tenth of that
budget, with a minimum of one second. This rule prevents one trial from
consuming the complete search, although it may be too restrictive for slow
datasets. Users can override it or disable evaluation timeouts explicitly.

Dataset-aware filtering is also heuristic. For classification, K-nearest
neighbors and Gaussian processes are removed at 10,000 samples; support-vector
classification and multilayer perceptrons are additionally removed at 50,000
samples. Similar rules are applied to regression, and multinomial naive Bayes
is removed when numeric features contain negative values. These rules reduce
obviously expensive or incompatible trials, but they also narrow the search and
therefore require ablation in Chapter 6.

\section{Current Limitations}

The following limitations are present in the code used for this draft:

\begin{itemize}[leftmargin=*]
  \item The search targets supervised tabular data and does not implement
        time-series, grouped, image, or text-specific validation.
  \item The surrogate uncertainty is estimated from variation among trees. It
        is a practical approximation, not a calibrated posterior distribution.
  \item The promotion quantile and dataset-size filtering thresholds are fixed
        heuristics rather than values established experimentally.
  \item Meta-learning behavior depends on the coverage and quality of the
        packaged historical metadata. A dataset unlike those records may gain
        little from warm starts.
  \item Process creation and serialization introduce overhead, particularly on
        Windows and for small, fast trials.
  \item Optional estimator libraries increase installation size and can behave
        differently across versions and hardware platforms.
  \item The focused automated tests verify specific integration paths but are
        not a substitute for benchmark comparisons against other AutoML tools.
\end{itemize}

These limitations are stated before the experimental chapters so that later
results can be interpreted against the actual boundaries of the implementation.

\section{User-Facing Example}

Listing~\ref{lst:usage} shows the intended public API.

\begin{lstlisting}[style=pythonstyle,caption={Classification with Accelera AutoML},label={lst:usage}]
from sklearn.datasets import make_classification
from sklearn.model_selection import train_test_split

from accelera.accelera_automl import AutoMLClassifier

X, y = make_classification(
    n_samples=1000,
    n_features=20,
    n_informative=12,
    random_state=42,
)
X_train, X_test, y_train, y_test = train_test_split(
    X, y, test_size=0.2, stratify=y, random_state=42
)

model = AutoMLClassifier(
    time_budget=300,
    n_trials=20,
    cv=3,
    random_state=42,
    search_n_parallel=2,
    n_jobs=1,
    use_meta_learning=True,
    use_ensemble=True,
    ensemble_strategy="stacked",
)
model.fit(X_train, y_train)

predictions = model.predict(X_test)
print(model.score(X_test, y_test))
print(model.return_leaderboard(top_n=3))
\end{lstlisting}

\section{Current Verification Assets}

Focused tests cover packaged meta-learning path resolution, compatibility of
the classification adapter with bagging, and regression stacked-ensemble
forward-selection state. Executable examples are provided for classification
and for regression using scikit-learn's diabetes dataset. These tests establish
basic integration behavior; the larger experimental protocol and quantitative
evaluation will be written in the reserved later chapters.

\section{Chapter Summary}

The implementation maps the proposed method to independently maintainable
packages. Public estimator classes manage user interaction, the engine
coordinates the workflow, evaluators own task-specific scoring, the optimizer
controls resource allocation, meta-learning provides initial configurations,
and ensemble modules construct the final combined candidates. This separation
supports future replacement or extension of individual subsystems without
changing the public estimator interface.

% Chapters 6--9 (Experimental Methodology, Results and Evaluation, Discussion,
% and Conclusion/Future Work) are intentionally reserved for a later revision.

\begin{thebibliography}{99}

\bibitem{autoweka}
Chris Thornton, Frank Hutter, Holger H. Hoos, and Kevin Leyton-Brown.
\newblock Auto-WEKA: Combined selection and hyperparameter optimization of
classification algorithms.
\newblock In \emph{Proceedings of the 19th ACM SIGKDD International Conference
on Knowledge Discovery and Data Mining}, 2013.

\bibitem{autosklearn}
Matthias Feurer, Aaron Klein, Katharina Eggensperger, Jost Springenberg,
Manuel Blum, and Frank Hutter.
\newblock Efficient and robust automated machine learning.
\newblock In \emph{Advances in Neural Information Processing Systems}, 2015.

\bibitem{tpot}
Randal S. Olson, Nathan Bartley, Ryan J. Urbanowicz, and Jason H. Moore.
\newblock Evaluation of a tree-based pipeline optimization tool for automating
data science.
\newblock In \emph{Applications of Evolutionary Computation}, 2016.

\bibitem{flaml}
Chi Wang, Qingyun Wu, Markus Weimer, and Erkang Zhu.
\newblock FLAML: A fast and lightweight AutoML library.
\newblock In \emph{Proceedings of Machine Learning and Systems}, 2021.

\bibitem{autogluon}
Nick Erickson et al.
\newblock AutoGluon-Tabular: Robust and accurate AutoML for structured data.
\newblock \emph{arXiv preprint arXiv:2003.06505}, 2020.

\bibitem{smac}
Frank Hutter, Holger H. Hoos, and Kevin Leyton-Brown.
\newblock Sequential model-based optimization for general algorithm
configuration.
\newblock In \emph{Learning and Intelligent Optimization}, 2011.

\bibitem{randomsearch}
James Bergstra and Yoshua Bengio.
\newblock Random search for hyper-parameter optimization.
\newblock \emph{Journal of Machine Learning Research}, 13:281--305, 2012.

\bibitem{stacking}
David H. Wolpert.
\newblock Stacked generalization.
\newblock \emph{Neural Networks}, 5(2):241--259, 1992.

\end{thebibliography}

\end{document}
